% (User input window)

\begin{conclusion}[Denominator rigidity of fixed weight layer coefficient generating function: unified closed form and second-order recursion]\label{con:xi-terminal-zm-fixed-weight-generating-function-rigidity}
Follow the notation of conclusion \ref{con:xi-terminal-zm-order4-rational} and conclusion \ref{con:xi-terminal-zm-2d-recurrence-amk}, and note
\[
a_{m,k}:=[y^k]\,Z_m(y),\qquad
A_k(t):=\sum_{m\ge 0}a_{m,k}t^m=[y^k]F(t,y),
\]
in
\[
F(t,y)=\sum_{m\ge 0}Z_m(y)t^m
=\frac{1+yt-yt^2-y^2t^3}{1-t-(2y+1)t^2+t^3+y(y+1)t^4}.
\]
Then for every fixed integer \(k\ge 0\), there is a unique polynomial \(\mathcal{P}_k(t)\in\ZZ[t]\) such that
\[
\boxed{\;
A_k(t)=\frac{\mathcal{P}_k(t)}{(1-t)^{2k+2}(1+t)^{k+1}}\;}
\]
And \(\mathcal{P}_k\) is uniquely determined by explicit second-order recursion: let
\[
\mathcal{P}_0(t)=1,\qquad
\mathcal{P}_1(t)=-t\bigl(t^4-t^3-1\bigr),\qquad
\mathcal{P}_2(t)=t^3\bigl(t^5-2t^4-t^3+t^2+t+1\bigr),
\]
Then for all \(k\ge 3\) we have
\[
\boxed{\;
\mathcal{P}_k(t)=(2t^2-t^4)\mathcal{P}_{k-1}(t)-t^4(1-t)^2(1+t)\mathcal{P}_{k-2}(t)\; }.
\]
\end{conclusion}
\begin{proof}[Proof sketch]
Write the denominator as
\[
D(t,y):=1-t-(2y+1)t^2+t^3+y(y+1)t^4=D_0(t)+yD_1(t)+y^2D_2(t),
\]
in
\[
D_0(t)=(1-t)^2(1+t),\qquad D_1(t)=-2t^2+t^4,\qquad D_2(t)=t^4,
\]
and write the numerator as
\(
N(t,y)=1+y(t-t^2)-y^2t^3
\)。
Write \(F(t,y)=\sum_{k\ge 0}A_k(t)y^k\), then by comparing the coefficients of \(y^k\) with the identity \(D(t,y)F(t,y)=N(t,y)\), we get
\[
D_0(t)A_k(t)+D_1(t)A_{k-1}(t)+D_2(t)A_{k-2}(t)=N_k(t),
\]
Among them \(N_0(t)=1\), \(N_1(t)=t-t^2\), \(N_2(t)=-t^3\), and \(N_k(t)=0\) (when \(k\ge 3\)), and it is agreed that \(A_{-1}=A_{-2}=0\).
Let \(\mathcal{P}_k(t):=D_0(t)^{k+1}A_k(t)\), then the above formula is equivalent to
\[
\mathcal{P}_k(t)=\bigl(2t^2-t^4\bigr)\mathcal{P}_{k-1}(t)-t^4D_0(t)\mathcal{P}_{k-2}(t)+N_k(t)D_0(t)^k.
\]
For \(k\ge 3\) when \(N_k\equiv 0\), we get the homogeneous second-order recursion shown;
The explicit initial value of \(k=0,1,2\) is obtained by directly simplifying the above formula.
Since \(D_0,D_1,D_2,N_k\in\ZZ[t]\) and the initial value is an integer coefficient, the recursion gives \(\mathcal{P}_k\in\ZZ[t]\).
Finally, the denominator structure is obtained by \(D_0(t)^{k+1}=(1-t)^{2k+2}(1+t)^{k+1}\). Certification completed.
\end{proof}

\begin{conclusion}[Quasipolynomiality of fixed weight layer coefficients: period \(2\), degree quantization and closed form of leading coefficient]\label{con:xi-terminal-zm-fixed-weight-quasipolynomial}
Follow the notation of conclusion \ref{con:xi-terminal-zm-order4-rational} and conclusion \ref{con:xi-terminal-zm-2d-recurrence-amk}, and note
\[
a_{m,k}:=[y^k]\,Z_m(y)\qquad(m\ge 0,\ k\ge 0).
\]
Then for each fixed integer \(k\ge 0\), there is a unique pair of rational coefficient polynomials \(P_k,Q_k\in\QQ[m]\), so that it is true for all \(m\ge 0\)
\[
\boxed{\;a_{m,k}=P_k(m)+(-1)^m Q_k(m)\;}
\]
And the number of times is satisfied
\[
\deg P_k=2k+1,\qquad \deg Q_k=k.
\]
Furthermore, the leading coefficients of \(P_k\) and \(Q_k\) have closed form
\[
\bigl[m^{2k+1}\bigr]P_k(m)=\frac{1}{2^{k+1}(2k+1)!},
\qquad
\bigl[m^{k}\bigr]Q_k(m)=\frac{1}{2^{2k+2}k!}.
\]
Equivalently, let \(m=2n+\varepsilon\)(\(\varepsilon\in\{0,1\}\)), then the two odd and even subsequences \(n\mapsto a_{2n,k}\) and \(n\mapsto a_{2n+1,k}\) are both \(n\) degree \(2k+1\) polynomials.
\end{conclusion}
\begin{proof}[Proof sketch]
From the conclusion \ref{con:xi-terminal-zm-fixed-weight-generating-function-rigidity},
There is an explicit closed form for every fixed \(k\ge 0\)
\[
A_k(t)=\frac{\mathcal{P}_k(t)}{(1-t)^{2k+2}(1+t)^{k+1}}\in\QQ(t),
\]
Thus \(A_k(t)\) only has poles at \(t=\pm 1\), and the pole orders at \(t=1\) and \(t=-1\) are exactly \(2k+2\) and \(k+1\) respectively.
Make partial fractions of \(A_k\) and take the coefficients item by item to get the decomposition that holds for all \(m\ge 0\)
\(
a_{m,k}=P_k(m)+(-1)^mQ_k(m)
\)
And the degree bound; the uniqueness comes from: if \(P(m)+(-1)^mQ(m)\equiv 0\) holds for all integers \(m\ge 0\), then there must be \(P\equiv 0\) and \(Q\equiv 0\).
\par
The leading coefficient is given by the highest-order pole at \(t=1\): by the recursion of the conclusion \ref{con:xi-terminal-zm-fixed-weight-generating-function-rigidity} substituting at \(t=1\) we can induce \(\mathcal{P}_k(1)=1\); and \((1+t)^{-(k+1)}\sim 2^{-(k+1)}\) (\(t\to 1\)), so
\(
A_k(t)\sim 2^{-(k+1)}(1-t)^{-(2k+2)}
\),thereby
\(
[t^m]A_k(t)=2^{-(k+1)}\frac{m^{2k+1}}{(2k+1)!}+O(m^{2k})
\)
gives the leading coefficient shown.
\par
For the leading coefficient of \(t=-1\), record the denominator
\[
D(t,y)=D_0(t)+yD_1(t)+y^2D_2(t),
\qquad
D_0(t)=D(t,0),\ \ D_1(t)=\partial_y D(t,0)=-2t^2+t^4,
\]
but
\(
D_0(t)\sim 4(1+t)
\)
(\(t\to-1\)) and \(D_1(-1)=-1\).
Since the numerator of \(F(t,y)\) takes the value \(1\) at \(y=0\), and its \(y\)-degree is at most \(2\), the highest-order pole of \(A_k(t)\) at \(t=-1\) is completely determined by \([y^k]\bigl(1/D(t,y)\bigr)\).
Write
\[
\frac{1}{D(t,y)}=\frac{1}{D_0(t)}
\left(1+\frac{yD_1(t)+y^2D_2(t)}{D_0(t)}\right)^{-1},
\]
Then the highest order term of \([y^k]\bigl(1/D(t,y)\bigr)\) is \((-1)^kD_1(t)^k/D_0(t)^{k+1}\);
Combining \(D_0(t)\sim 4(1+t)\) and \(D_1(-1)=-1\) we get:
\[
A_k(t)=\frac{1}{4^{k+1}}(1+t)^{-(k+1)}+O\!\bigl((1+t)^{-k}\bigr)\qquad(t\to-1),
\]
Therefore, the pole order of \(t=-1\) is exactly \(k+1\), so \(\deg Q_k=k\).
and by
\(
[t^m](1+t)^{-(k+1)}=(-1)^m\binom{m+k}{k}
\)
get
\(
[m^k]Q_k(m)=4^{-(k+1)}/k!=2^{-(2k+2)}/k!
\)。
The above decomposition and subsequent low-order closure can be reviewed by the script \texttt{scripts/exp\_fold\_zm\_low\_weight\_quasipolynomial\_audit.py}. Certification completed.
\end{proof}

\begin{conclusion}[Explicitization of two-term principal parts: polynomial principal growth and order separation of alternating oscillations]\label{con:xi-terminal-zm-fixed-weight-two-term-asymptotic}
Follow the notation of the conclusion \ref{con:xi-terminal-zm-fixed-weight-generating-function-rigidity} and note \(\mathcal{P}_k(t)\in\ZZ[t]\) as the numerator of \(A_k(t)\).
Then for all \(k\ge 0\) we have
\[
\boxed{\;
\mathcal{P}_k(1)=\mathcal{P}_k(-1)=1,\qquad \mathcal{P}_k'(1)=0\;}
\]
And when \(k\ge 1\), the first law of the second and third order derivatives is established:
\[
\boxed{\;
\mathcal{P}_k''(1)=4-12k,\qquad \mathcal{P}_k^{(3)}(1)=42-78k\; }.
\]
Therefore, for every fixed \(k\) (let \(m\to\infty\)) there are two explicit principal expansions
\[
\boxed{\;
a_{m,k}
=\frac{1}{2^{k+1}}\binom{m+2k+1}{2k+1}
\frac{k+1}{2^{k+2}}\binom{m+2k}{2k}
O_k(m^{2k-1})
(-1)^m O_k(m^{k})\; }.
\]
where \(O_k(m^{2k-1})\) comes from the next higher order pole of \(t=1\), and the alternating term \((-1)^mO_k(m^k)\) comes from the pole of \(t=-1\).
\end{conclusion}
\begin{proof}[Proof sketch]
Recursion from the conclusion \ref{con:xi-terminal-zm-fixed-weight-generating-function-rigidity}
\[
\mathcal{P}_k(t)=(2t^2-t^4)\mathcal{P}_{k-1}(t)-t^4(1-t)^2(1+t)\mathcal{P}_{k-2}(t)
\]
Substituting at \(t=\pm 1\) we get \(\mathcal{P}_k(\pm1)=\mathcal{P}_{k-1}(\pm1)=1\).
Taking the derivative of this recurrence at \(t=1\) and noting that \((2t^2-t^4)'|_{t=1}=0\) and \(t^4(1-t)^2(1+t)\) vanishes to at least second order at \(t=1\), we get \(\mathcal{P}_k'(1)=\mathcal{P}_{k-1}'(1)=0\).
The laws of second and third derivatives are the same by induction by differentiating the recursion second and third at \(t=1\) (starting with a direct check of \(k=1,2\)).
\par
Finally by
\[
A_k(t)=\frac{\mathcal{P}_k(t)}{(1-t)^{2k+2}(1+t)^{k+1}}
\]
The local expansion at \(t=1\) and \(t=-1\) operates as coefficient extraction:
At \(t=1\),
\(
(1+t)^{-(k+1)}=2^{-(k+1)}\bigl(1+\frac{k+1}{2}(1-t)+O((1-t)^2)\bigr)
\)
And combine \(\mathcal{P}_k(1)=1\), \(\mathcal{P}_k'(1)=0\) to get the two main parts shown;
And at \(t=-1\), the order of the alternating term obtained from the pole order \(k+1\) is \(O_k(m^k)\).
Certification completed.
\end{proof}

\begin{conclusion}[The total closed form of the low-weight layer coefficient: \(k\le 3\) holds for all \(m\ge 0\) point by point]\label{con:xi-terminal-zm-fixed-weight-kle3-closed}
In the notation of the conclusion \ref{con:xi-terminal-zm-fixed-weight-quasipolynomial}, let \(n\ge 0\).
Then the two odd-even subsequences \(\{a_{2n,k}\}_{n\ge 0}\) and \(\{a_{2n+1,k}\}_{n\ge 0}\) of \(k=0,1,2,3\) have the following point-wise closed formulas:
\[
a_{2n,0}=n+1,\qquad a_{2n+1,0}=n+1.
\]
\[
a_{2n,1}=\frac{1}{3}n^3+\frac{3}{2}n^2+\frac{1}{6}n,
\qquad
a_{2n+1,1}=\frac{1}{3}n^3+2n^2+\frac{5}{3}n+1.
\]
\[
a_{2n,2}=\frac{1}{30}n^5+\frac{3}{8}n^4-\frac{1}{12}n^3-\frac{7}{8}n^2+\frac{11}{20}n,
\]
\[
a_{2n+1,2}=\frac{1}{30}n^5+\frac{11}{24}n^4+\frac{3}{4}n^3-\frac{11}{24}n^2+\frac{13}{60}n.
\]
\[
a_{2n,3}=\frac{1}{630}n^7+\frac{1}{30}n^6+\frac{7}{360}n^5-\frac{17}{24}n^4+\frac{523}{360}n^3-\frac{33}{40}n^2+\frac{11}{420}n,
\]
\[
a_{2n+1,3}=\frac{1}{630}n^7+\frac{7}{180}n^6+\frac{23}{180}n^5-\frac{19}{36}n^4+\frac{29}{180}n^3+\frac{22}{45}n^2-\frac{61}{210}n.
\]
\end{conclusion}
\begin{proof}[Proof sketch]
For \(k\le 3\), it can be known from \(A_k(t)=[y^k]F(t,y)\in\QQ(t)\) and its pole is only located at \(t=\pm1\)
\(\{a_{2n,k}\}_{n\ge 0}\)、\(\{a_{2n+1,k}\}_{n\ge 0}\)
Both are polynomial sequences with respect to \(n\) (with degree up to \(2k+1\)).
Therefore, the unique polynomial can be recovered with finite point values ​​and checked point by point with the two-dimensional recursion (conclusion \ref{con:xi-terminal-zm-2d-recurrence-amk});
You can also directly perform partial fractions on \(A_k(t)\) and recover the coefficients.
The above closed form and its point-by-point consistency can be checked by the script \texttt{scripts/exp\_fold\_zm\_low\_weight\_quasipolynomial\_audit.py}. Certification completed.
\end{proof}

